% XaTanThuanDong - Project Report (single-file)
% Compile: xelatex main.tex  (recommended) or pdflatex
\documentclass[12pt,a4paper]{article}

\usepackage[a4paper,margin=2.2cm]{geometry}
\usepackage{fontspec}
\usepackage[vietnamese]{babel}
\usepackage{hyperref}
\usepackage{graphicx}
\usepackage{longtable}
\usepackage{booktabs}
\usepackage{array}
\usepackage{enumitem}
\usepackage{fancyhdr}
\usepackage{titlesec}
\usepackage{xcolor}
\usepackage{listings}
\usepackage{caption}

\setmainfont{Times New Roman}

\definecolor{codebg}{RGB}{248,248,248}
\definecolor{codeborder}{RGB}{220,220,220}
\lstdefinestyle{appcode}{
  backgroundcolor=\color{codebg},
  frame=single,
  rulecolor=\color{codeborder},
  basicstyle=\ttfamily\small,
  breaklines=true,
  columns=fullflexible,
  keepspaces=true,
}

\hypersetup{
  colorlinks=true,
  linkcolor=blue,
  urlcolor=blue,
  citecolor=blue,
  pdftitle={Báo cáo đồ án - XaTanThuanDong},
}

\titleformat{\section}{\Large\bfseries}{\thesection.}{0.6em}{}
\titleformat{\subsection}{\large\bfseries}{\thesubsection.}{0.6em}{}

\pagestyle{fancy}
\fancyhf{}
\lhead{Xa Tân Thuận Đông}
\rhead{Báo cáo dự án}
\cfoot{\thepage}

\begin{document}

% -------------------- Cover --------------------
\begin{titlepage}
  \centering
  \vspace*{1.5cm}
  {\Large \textbf{BÁO CÁO ĐỒ ÁN}}\\[0.4cm]
  {\LARGE \textbf{Website Xã Tân Thuận Đông}}\\[0.8cm]

  \vspace{0.4cm}
  \begin{tabular}{>{\bfseries}p{5cm}p{9cm}}
    Sinh viên thực hiện: & Tran Nguyen Phat Dat \\ 
    MSSV: & 0023411535 \\ 
    Dự án: & XaTanThuanDong (Frontend + Backend) \\ 
    Công nghệ: & React (Vite), .NET 8 Web API, EF Core \\ 
    Ngày báo cáo: & 24/03/2026 \\
  \end{tabular}

  \vfill
  {\large \textit{Tài liệu được tổng hợp từ cấu trúc source trong workspace:}\\
  \texttt{d:\\Tran Nguyen Phat Dat-0023411535\\github\\XaTanThuanDong}}

  \vspace*{1.2cm}
\end{titlepage}

\tableofcontents
\newpage

% -------------------- 1. Overview --------------------
\section{Tổng quan}
Dự án \textbf{XaTanThuanDong} là một hệ thống website cung cấp thông tin và dịch vụ hành chính công cho \textbf{Xã Tân Thuận Đông}. Hệ thống gồm:
\begin{itemize}[leftmargin=1.2cm]
  \item \textbf{Backend}: ASP.NET Core 8 Web API, Entity Framework Core, xác thực JWT, phân quyền Admin.
  \item \textbf{Frontend}: React + Vite, giao diện người dùng (public) và khu vực quản trị (admin).
\end{itemize}

\subsection{Mục tiêu}
\begin{itemize}[leftmargin=1.2cm]
  \item Cung cấp tin tức, thông báo, nội dung giới thiệu địa phương.
  \item Cung cấp danh sách dịch vụ, thủ tục, và tiếp nhận hồ sơ đăng ký dịch vụ (service application).
  \item Khu vực quản trị quản lý: người dùng, nội dung (bài viết, danh mục), dịch vụ, dashboard thống kê.
\end{itemize}

% -------------------- 2. Structure --------------------
\section{Cấu trúc thư mục (tóm tắt)}
\subsection{Root}
\begin{lstlisting}[style=appcode]
XaTanThuanDong.sln
backend/
  XaTanThuanDong.Api/
frontend/
  xatanthuandong-web/
\end{lstlisting}

\subsection{Backend: \texttt{backend/XaTanThuanDong.Api}}
Có các thành phần chính: Controllers, Data (DbContext + Seed), Dtos, Models, Services, Validators, Middlewares, Migrations, Uploads.

Controllers quan trọng:
\begin{itemize}[leftmargin=1.2cm]
  \item \texttt{AuthController.cs}: đăng nhập/đăng ký, cấp JWT.
  \item \texttt{PublicContentController.cs}: API public cho bài viết, danh mục.
  \item \texttt{PublicServicesController.cs}: API public cho dịch vụ, thủ tục.
  \item \texttt{Admin*Controller.cs}: API quản trị (dashboard, users, services, content).
\end{itemize}

\subsection{Frontend: \texttt{frontend/xatanthuandong-web}}
Tổ chức theo standard Vite + React:
\begin{itemize}[leftmargin=1.2cm]
  \item \texttt{src/pages}: các trang public (Home, News list/detail, Services, About, Contact, Gallery) và admin.
  \item \texttt{src/lib}: helper gọi API và auth.
  \item \texttt{vite.config.js}: cấu hình dev server (proxy API trong dev).
\end{itemize}

% -------------------- 3. Architecture --------------------
\section{Kiến trúc hệ thống}
\subsection{Luồng tổng quan}
\begin{enumerate}[leftmargin=1.2cm]
  \item Người dùng truy cập frontend (React).
  \item Frontend gọi các endpoint public để lấy tin tức/dịch vụ.
  \item Admin đăng nhập, nhận JWT, sau đó gọi các endpoint admin với header Authorization.
  \item Backend truy cập CSDL qua EF Core (DbContext), trả JSON.
\end{enumerate}

\subsection{Công nghệ}
\begin{longtable}{p{5.2cm}p{9.2cm}}
\toprule
\textbf{Hạng mục} & \textbf{Công nghệ} \\
\midrule
Backend & ASP.NET Core 8 Web API, EF Core Migrations, FluentValidation, JWT, Middleware xử lý lỗi \\
Frontend & React, Vite, React Router, Axios (thông qua \texttt{src/lib/api.js}) \\
CSDL & (Theo cấu hình dự án) EF Core, Migration \texttt{InitialCreate} \\
\bottomrule
\end{longtable}

% -------------------- 4. Database --------------------
\section{Cơ sở dữ liệu (mô hình thực thể)}
Dựa theo thư mục \texttt{Models/} trong backend, các thực thể chính gồm:
\begin{itemize}[leftmargin=1.2cm]
  \item \textbf{AppUser}: người dùng hệ thống (admin/user).
  \item \textbf{Category}: danh mục tin.
  \item \textbf{Article}: bài viết/tin tức.
  \item \textbf{Comment}: bình luận (nếu sử dụng).
  \item \textbf{Media}: tài nguyên media (ảnh/file).
  \item \textbf{ServiceProcedure}: thủ tục dịch vụ.
  \item \textbf{ServiceApplication}: hồ sơ đăng ký/đơn yêu cầu dịch vụ.
  \item \textbf{AuditLog}: log thao tác (nếu dùng).
\end{itemize}

\subsection{Bảng tóm tắt entities}
\begin{longtable}{p{4.2cm}p{10.2cm}}
\toprule
\textbf{Entity} & \textbf{Mô tả} \\
\midrule
AppUser & Lưu thông tin tài khoản, role, password hash, ... \\
Category & Danh mục nội dung/bài viết \\
Article & Nội dung tin tức, slug, title, summary, body, thumbnail, publishedAt \\
ServiceProcedure & Thông tin thủ tục: tên, mô tả, hồ sơ cần, thời gian, lệ phí,... \\
ServiceApplication & Yêu cầu dịch vụ của người dân: thông tin liên hệ, trạng thái xử lý, đính kèm,... \\
Media & Quản lý file/ảnh upload \\
AuditLog & Ghi nhận lịch sử thao tác/quản trị \\
Comment & Bình luận theo bài viết \\
\bottomrule
\end{longtable}

\textit{Ghi chú: Chi tiết cột (columns), quan hệ (FK) và index được định nghĩa trong Migration: \texttt{Migrations/20260323004718\_InitialCreate.cs}.}

% -------------------- 5. API --------------------
\section{API (tóm tắt endpoint)}
\subsection{Public API}
Các endpoint public nằm ở \texttt{PublicContentController.cs} và \texttt{PublicServicesController.cs}. Ví dụ:
\begin{itemize}[leftmargin=1.2cm]
  \item \texttt{GET /api/public/categories} --- lấy danh mục.
  \item \texttt{GET /api/public/articles} --- danh sách bài viết (hỗ trợ query như \texttt{q}, \texttt{categoryId}).
  \item \texttt{GET /api/public/articles/\{slug\}} --- chi tiết bài viết.
  \item \texttt{GET /api/public/services} --- danh sách dịch vụ/thủ tục.
\end{itemize}

\subsection{Auth API}
\texttt{AuthController.cs} cung cấp luồng đăng nhập/đăng ký và cấp token JWT.

\subsection{Admin API}
Các endpoint quản trị (yêu cầu JWT + role phù hợp):
\begin{itemize}[leftmargin=1.2cm]
  \item \texttt{/api/admin/dashboard} --- số liệu thống kê.
  \item \texttt{/api/admin/users} --- quản lý user.
  \item \texttt{/api/admin/content} --- quản lý categories/articles/media.
  \item \texttt{/api/admin/services} --- quản lý dịch vụ/thủ tục/hồ sơ.
\end{itemize}

% -------------------- 6. Frontend pages --------------------
\section{Frontend (các trang chính)}
Các trang nằm ở \texttt{frontend/xatanthuandong-web/src/pages}.

\begin{longtable}{p{5.2cm}p{9.2cm}}
\toprule
\textbf{Trang} & \textbf{Mô tả}\\
\midrule
HomePage.jsx & Trang chủ, giới thiệu nhanh và các block nội dung \\
NewsListPage.jsx & Danh sách tin tức, tìm kiếm theo từ khóa/danh mục \\
NewsDetailPage.jsx & Chi tiết bài viết theo slug \\
ServicesPage.jsx & Danh sách dịch vụ/thủ tục \\
AboutPage.jsx & Giới thiệu \\
ContactPage.jsx & Liên hệ \\
GalleryPage.jsx & Thư viện hình ảnh \\
Admin/* & Khu vực quản trị (tùy triển khai) \\
\bottomrule
\end{longtable}

% -------------------- 7. Run instructions --------------------
\section{Hướng dẫn chạy dự án (Development)}
\subsection{Yêu cầu}
\begin{itemize}[leftmargin=1.2cm]
  \item .NET SDK 8.x
  \item Node.js (khuyến nghị LTS)
\end{itemize}

\subsection{Chạy Backend}
Mở solution \texttt{XaTanThuanDong.sln} hoặc chạy project API.

Chạy API (ví dụ):
\begin{lstlisting}[style=appcode]
cd backend/XaTanThuanDong.Api
# dotnet run
\end{lstlisting}

Cấu hình port xem trong \texttt{Properties/launchSettings.json}. Mặc định backend chạy: \texttt{http://localhost:5168} (phụ thuộc launch settings).

\subsection{Chạy Frontend}
\begin{lstlisting}[style=appcode]
cd frontend/xatanthuandong-web
# npm install
# npm run dev
\end{lstlisting}

Trong dev, frontend nên gọi API theo dạng tương đối \texttt{/api/...} và cấu hình proxy trong \texttt{vite.config.js} để trỏ về backend.

% -------------------- 8. Notes --------------------
\section{Ghi chú triển khai và bảo mật}
\begin{itemize}[leftmargin=1.2cm]
  \item Không commit secret (JWT secret, connection string production).
  \item Bật HTTPS, cấu hình CORS đúng domain production.
  \item Với upload: giới hạn loại file, giới hạn dung lượng, đặt tên file an toàn.
  \item Ghi log và xử lý exception tập trung qua middleware.
\end{itemize}

\section{Kết luận}
Tài liệu này gom tóm tắt cấu trúc và chức năng chính của hệ thống. Có thể mở rộng phần báo cáo bằng cách bổ sung: ERD chi tiết, ảnh chụp màn hình UI, test case, và kết quả triển khai.

\end{document}
